%
% 16-quaternionen.tex -- Quaternionen
%
% (c) 2026 Prof Dr Andreas Müller
%
\subsection{Quaternionen}
Die Matrix $J$ ist nicht die einzige Matrix mit der Eigenschaft
$J^2=-I$.
\index{Quaternionen}%
Lässt man komplexe Einträge in den Matrizen zu, haben auch die
Matrizen
\begin{align*}
-i\sigma_x
&=
-i
\begin{pmatrix*}[r]
0&1\\
1&0
\end{pmatrix*}
=
\begin{pmatrix*}[r]
 0&-i\\
-i& 0
\end{pmatrix*}
&&\Rightarrow&
(-i\sigma_x)^2
&=
\begin{pmatrix*}[r]
-1& 0\\
 0&-1
\end{pmatrix*}
=
-I
\\
-i\sigma_y
&=
-i
\begin{pmatrix*}[r]
0&-i\\
i& 0
\end{pmatrix*}
=
\begin{pmatrix*}[r]
0&-1\\
1&0
\end{pmatrix*}
=
J
&&\Rightarrow&
(-i\sigma_y)^2
&=
J^2
=
-I
\\
-i\sigma_z
&=
-i
\begin{pmatrix*}[r]
1&0\\
0&-1
\end{pmatrix*}
=
\begin{pmatrix*}[r]
-i&0\\
 0&i
\end{pmatrix*}
&&\Rightarrow&
(-i\sigma_z)^2
&=
\begin{pmatrix*}[r]
-1& 0\\
 0&-1
\end{pmatrix*}
=
-I.
\end{align*}
diese Eigenschaft.

\begin{definition}[Pauli-Matrizen]
Die Matrizen
\begin{align*}
\sigma_x
&=
\begin{pmatrix*}[r]
0&1\\
1&0
\end{pmatrix*},
&
\sigma_y
&=
\begin{pmatrix*}[r]
0&-i\\
i& 0
\end{pmatrix*}
&&\text{und}
&
\sigma_z
&=
\begin{pmatrix*}[r]
1&0\\
0&-1
\end{pmatrix*}
\end{align*}
heissen die \emph{Pauli-Matrizen}.
\index{Pauli-Matrix}%
\end{definition}

Die Pauli-Matrizen spielen eine wichtige Rolle in der Beschreibung des
Spins in der Quantenmechanik.
\index{Quantenmechanik}%
\index{Spin}%
Sie können auch verwendet werden, um die Gruppe $\operatorname{SU}(2)$ 
der spezielle unitären Matrizen zu beschreiben.
\index{unitär}%

Schreiben wir
\begin{align}
\boldsymbol{i}
&=
-i\sigma_x
=
\begin{pmatrix}
0&-i\\
-i&0
\end{pmatrix}
,
&
\boldsymbol{j}
&=
-i\sigma_y
=
\begin{pmatrix}
0&-1\\
1&0
\end{pmatrix}
=
J
&&\text{und}
&
\boldsymbol{k}
&=
-i\sigma_z
=
\begin{pmatrix}
-i & 0 \\
 0 & i
\end{pmatrix}
,
\label{buch:komplex:zahlen:eqn:quaternionen}
\end{align}
erhalten wir zusätzlich zu
$\boldsymbol{i}^2
=
\boldsymbol{j}^2
=
\boldsymbol{k}^2
=-1$
durch Nachrechnen die folgenden Beziehungen:
\begin{align}
\boldsymbol{i}\boldsymbol{j}
&=
\begin{pmatrix*}[r]
 0&-i\\
-i& 0
\end{pmatrix*}
\begin{pmatrix*}[r]
 0&-1\\
 1& 0
\end{pmatrix*}
=
\begin{pmatrix*}[r]
-i & 0 \\
 0 & i
\end{pmatrix*}
=
\boldsymbol{k},
&
\boldsymbol{j}\boldsymbol{i}
&=
\begin{pmatrix*}[r]
 0&-1\\
 1& 0
\end{pmatrix*}
\begin{pmatrix*}[r]
 0&-i\\
-i& 0
\end{pmatrix*}
=
\begin{pmatrix*}[r]
i & 0 \\
0 & -i
\end{pmatrix*}
=
-\boldsymbol{k},
\label{buch:komplex:zahlen:eqn:quatrel}
\\
\boldsymbol{j}\boldsymbol{k}
&=
\begin{pmatrix*}[r]
 0&-1\\
 1& 0
\end{pmatrix*}
\begin{pmatrix*}[r]
-i & 0 \\
 0 & i
\end{pmatrix*}
=
\begin{pmatrix*}[r]
 0&-i\\
-i& 0
\end{pmatrix*}
=
\boldsymbol{i},
&
\boldsymbol{k}\boldsymbol{j}
&=
\begin{pmatrix*}[r]
-i & 0 \\
 0 & i
\end{pmatrix*}
\begin{pmatrix*}[r]
 0&-1\\
 1& 0
\end{pmatrix*}
=
\begin{pmatrix*}[r]
0&i\\
i&0
\end{pmatrix*}
=
-\boldsymbol{i},
\notag
\\
\boldsymbol{k}\boldsymbol{i}
&=
\begin{pmatrix*}[r]
-i & 0 \\
 0 & i
\end{pmatrix*}
\begin{pmatrix*}[r]
 0&-i\\
-i& 0
\end{pmatrix*}
=
\begin{pmatrix*}[r]
0 & -1 \\
1 & 0
\end{pmatrix*}
=
\boldsymbol{j},
&
\boldsymbol{i}\boldsymbol{k}
&=
\begin{pmatrix*}[r]
 0&-i\\
-i& 0
\end{pmatrix*}
\begin{pmatrix*}[r]
-i & 0 \\
 0 & i
\end{pmatrix*}
=
\begin{pmatrix*}[r]
 0 & 1 \\
-1 & 0
\end{pmatrix*}
=
-\boldsymbol{j}.
\notag
\end{align}
Wir zu erwarten ist, sind die Produkte der Matrizen
$\boldsymbol{i}$,
$\boldsymbol{j}$
und
$\boldsymbol{k}$
nicht kommutativ.
Diese Relationen können auch in die Gleichungen
\begin{equation}
\boldsymbol{i}^2
=
\boldsymbol{j}^2
=
\boldsymbol{k}^2
=
\boldsymbol{i}
\boldsymbol{j}
\boldsymbol{k}
=
-1
\label{buch:komplex:zahlen:eqn:hamiltonrelationen}
\end{equation}
zusammengefasst werden.
Multipliziert man die letzte Identität von rechts mit $\boldsymbol{k}$,
erhält man
\input{chapters/010-komplex/fig/fig-hamilton.tex}%
\begin{equation*}
-\boldsymbol{k}
=
(
\boldsymbol{i}
\boldsymbol{j}
\boldsymbol{k}
)
\boldsymbol{k}
=
\boldsymbol{i}
\boldsymbol{j}
\boldsymbol{k}^2
=
-
\boldsymbol{i}
\boldsymbol{j}
\quad\Rightarrow\quad
\boldsymbol{i}
\boldsymbol{j}
=
\boldsymbol{k},
\end{equation*}
also die erste Identität in \eqref{buch:komplex:zahlen:eqn:quatrel}.
Die anderen Identitäten können auf ähnliche Weise aus
\eqref{buch:komplex:zahlen:eqn:hamiltonrelationen}
gewonnen werden.

Die Menge
\[
\mathbb{H}
=
\langle 1,
\boldsymbol{i},
\boldsymbol{j},
\boldsymbol{k}
\rangle
=
\{
x_0
+
x_1\boldsymbol{i}
+
x_2\boldsymbol{j}
+
x_3\boldsymbol{k}
\mid
x_0,\dots,x_3\in\mathbb{R}
\}
\subset
M_{2\times 2}(\mathbb{C})
\]
ist einerseits ein reeller Vektorraum und andererseits eine
Algebra mit den Rechenregeln
\eqref{buch:komplex:zahlen:eqn:hamiltonrelationen},
die William Rowan Hamilton im Jahr 1843 entdeckt hat.
Sie heisst die Algebra der \emph{Quaternionen}.
Hamilton sollen die algebraischen Regeln
\index{Hamilton, William Rowan}%
\eqref{buch:komplex:zahlen:eqn:hamiltonrelationen}
während eines Spaziergangs über die Broom-Bridge in Dublin
\index{Dublin}%
\index{Broom-Bridge}%
eingefallen sein und er soll sie spontan in einen Stein
geritzt haben, woran die in Abbildung~\ref{buch:komplex:fig:hamilton}
abgebildete, an der Broom-Bridge angebrachte Plakette erinnert.

Die Matrizen $\boldsymbol{i}$ und $\boldsymbol{k}$ enthalten
Einträge, die selbst komplexe Zahlen sind.
Um eine Darstellung ganz ohne komplexe Zahlen zu bekommen, muss
jeder Eintrag der Matrizen in \eqref{buch:komplex:zahlen:eqn:quaternionen}
durch eine $2\times 2$-Matrix ersetzt werden.
Dadurch entstehen drei $4\times 4$-Matrizen, die wir mit grossen
fetten Buchstaben schreiben:
\begin{align*}
\boldsymbol{I}
&=
\begin{pmatrix*}[r]
 0& 0& 0& 1\\
 0& 0&-1& 0\\
 0& 1& 0& 0\\
-1& 0& 0& 0
\end{pmatrix*},
&
\boldsymbol{J}
&=
\begin{pmatrix*}[r]
 0& 0&-1& 0\\
 0& 0& 0&-1\\
 1& 0& 0& 0\\
 0& 1& 0& 0
\end{pmatrix*}
&&\text{und}
&
\boldsymbol{K}
&=
\begin{pmatrix*}[r]
 0& 1& 0& 0\\
-1& 0& 0& 0\\
 0& 0& 0&-1\\
 0& 0& 1& 0
\end{pmatrix*}.
\end{align*}
Man kann nachrechnen, dass die für
$\boldsymbol{i}$,
$\boldsymbol{j}$
und
$\boldsymbol{k}$
gefundenen Rechenregeln
\eqref{buch:komplex:zahlen:eqn:hamiltonrelationen}
auch für
$\boldsymbol{I}$,
$\boldsymbol{J}$
und
$\boldsymbol{K}$
gelten.

So wie Drehungen der komplexen Ebene durch Multiplikation mit
komplexen Zahlen beschrieben werden, können Quaternionen Drehungen
des dreidmensionalen Raumes beschreiben.
Dies wird im Kapitel~\ref{chapter:qa} beschrieben.


